\documentclass[11pt]{article}

\usepackage[a4paper, margin=1in]{geometry}
\usepackage{times}
\usepackage{setspace}
\usepackage{natbib}
\usepackage{hyperref}
\usepackage{graphicx}
\usepackage{booktabs}
\usepackage{authblk}

\doublespacing
\setlength{\parskip}{0.5em}

\title{External longitudinal validation of metabolic syndrome risk scores for cardiovascular, diabetes, and all-cause mortality in the US NHANES Linked Mortality File}

\author[1]{Jay Hemnani}
\affil[1]{Independent researcher. Correspondence: jayhemnani992000@gmail.com.}

\date{}

\begin{document}

\maketitle

\begin{abstract}
\noindent \textbf{Importance.} Two metabolic syndrome risk scoring methods, a continuous score based on triangular areal similarity (RMRS) and a decision-tree formulation, have been developed and internally validated on Korean cohorts. External validation in a US cohort against established clinical risk equations, with long-term mortality outcomes, has not been reported.

\noindent \textbf{Objective.} To externally validate the RMRS and the metabolic syndrome decision tree against the ACC/AHA Pooled Cohort Equations (PCE), Framingham 2008, and FINDRISC for cardiovascular, diabetes-related, and all-cause mortality in a representative US adult population.

\noindent \textbf{Design, setting, and participants.} Prospective cohort using NHANES survey cycles 1999 through 2018 linked to the public-use NHANES Linked Mortality File (2019 release). Adults aged 40 to 79 years with non-missing data on the fasting metabolic syndrome components were eligible. Cohort N and exclusions reported in the Results.

\noindent \textbf{Exposures.} Five risk scores computed at NHANES baseline: RMRS, the metabolic syndrome decision tree, PCE, Framingham 2008, and FINDRISC.

\noindent \textbf{Main outcomes and measures.} Cardiovascular mortality, diabetes-related mortality (broadened definition; see Methods), and all-cause mortality, ascertained via the Linked Mortality File. Discrimination was quantified by inverse-probability-of-censoring weighted (IPCW) time-dependent AUC at 5 years and at off-cap horizons (9.5 years for cardiovascular and all-cause; 14.5 years for diabetes-related). Calibration, net benefit (decision curve analysis), pairwise delta-AUC, and incremental value (continuous NRI and IDI) were reported per registered analysis plan. Confidence intervals were computed via PSU-cluster bootstrap at 500 resamples.

\noindent \textbf{Results.} The pooled fasting-subsample cohort comprised 17{,}031 adults, with 806 all-cause deaths and 181 cardiovascular deaths accrued within the 10-year follow-up cap and 80 diabetes-related deaths under the broadened definition within the 15-year cap. For cardiovascular and all-cause mortality at the 9.5-year horizon, the clinical risk equations clearly dominated the metabolic-syndrome-derived scores; Framingham achieved a bootstrap-estimated AUC of 0.859 (95\% CI 0.833 to 0.883) for cardiovascular mortality and 0.810 (0.795 to 0.825) for all-cause mortality, compared with RMRS at 0.662 (0.625 to 0.692) and 0.595 (0.574 to 0.618). For diabetes-related mortality at the 14.5-year horizon, RMRS and FINDRISC were comparable (RMRS 0.747 [0.692 to 0.800] vs FINDRISC 0.766 [0.715 to 0.817]), and both clearly exceeded the metabolic syndrome decision tree at 0.585 (0.533 to 0.633). The registered H3 non-inferiority test (RMRS vs FINDRISC for diabetes-related mortality, margin $-0.05$ AUC) yielded a point delta-AUC of $+0.011$ favoring RMRS with a bootstrap 95\% CI of $(-0.053, +0.075)$, in which the lower bound exceeded the registered margin by 0.003 AUC units. The DeLong-based CI agreed to a thousandth. Adding RMRS on top of FINDRISC for diabetes-related mortality produced an IDI of 0.0031 (0.0002 to 0.0099) and a continuous NRI of 0.173 (0.033 to 0.287). An XGBoost ceiling on the same five components placed an upper bound of 0.799 AUC for diabetes-related mortality at the 14.5-year horizon, about 5 AUC points above RMRS and FINDRISC.

\noindent \textbf{Conclusions and relevance.} Metabolic-syndrome-derived risk scores do not substitute for established clinical risk equations when the outcome is cardiovascular or all-cause mortality. For diabetes-related mortality, RMRS matches the dedicated FINDRISC screening tool on point discrimination, although the registered non-inferiority margin was not formally established by the bootstrap CI. The decision-tree formulation underperformed the continuous RMRS on every outcome, suggesting that the underperformance pattern reflects the tree structure rather than the metabolic syndrome signal.
\end{abstract}

\section*{Introduction}

Metabolic syndrome is a clustering of central adiposity, dysglycemia, dyslipidemia, and elevated blood pressure that is associated with increased risk of cardiovascular disease, type 2 diabetes, and all-cause mortality \citep{ncepatpiii, idfmets}. Several scoring frameworks have been proposed to quantify metabolic syndrome severity on a continuous scale or via decision-tree classification, with the aim of improving prevention beyond the binary NCEP ATP III or IDF criteria.

Two such frameworks have recently been developed and internally validated on Korean cohorts. The Robust Metabolic Syndrome Risk Score (RMRS) maps the five component features into a continuous risk index via triangular areal similarity \citep{shin2024rmrs}. A complementary decision-tree formulation segments the input space using the same five component features \citep{shin2023b9}. The two papers together constitute the most-cited line of research in this group's metabolic syndrome work, with the decision-tree paper being the highest-cited.

External validation of these scores in a non-Asian population, against established clinical prediction equations, on long-term mortality outcomes, has not been reported. The ACC/AHA Pooled Cohort Equations \citep{goff2014pce} and the Framingham 2008 general cardiovascular risk profile \citep{dagostino2008framingham} are the primary baselines for cardiovascular risk prediction in US clinical practice. FINDRISC \citep{lindstrom2003findrisc} is the most widely used non-invasive screening tool for incident type 2 diabetes. A head-to-head longitudinal comparison on a representative US cohort would clarify the place, if any, of metabolic-syndrome-derived scores within the existing prediction-tool ecosystem.

We performed this comparison using NHANES survey cycles 1999 through 2018 linked to the public-use Linked Mortality File (2019 release) \citep{nhaneslmf2019}. The analysis was pre-registered on OSF, follows TRIPOD+AI \citep{tripodai2024} and STROBE \citep{strobe2007} reporting standards, and reports discrimination, calibration, net benefit, pairwise delta-AUC, and incremental value for three mortality outcomes (cardiovascular, diabetes-related, all-cause) across all five risk scores. The registered primary non-inferiority hypothesis (H3) tested whether RMRS is non-inferior to FINDRISC for predicting diabetes-related mortality at the 14.5-year horizon, with a registered margin of $-0.05$ AUC.

\section*{Methods}

\subsection*{Study design and data sources}

We used a prospective cohort design on the public-use NHANES survey \citep{nhaneslmf2019} cycles A through J (calendar years 1999 through 2018) linked to the National Center for Health Statistics public-use Linked Mortality File (LMF) 2019 release. Both files are de-identified, publicly available, and do not require IRB approval for secondary analysis. NHANES employs a complex multi-stage probability sample of the US civilian, non-institutionalized population with sampling stratification (\texttt{sdmvstra}), primary sampling units (\texttt{sdmvpsu}), and per-cycle survey weights. All analyses incorporated the survey design via the R \texttt{survey} package; effective survey weights were rescaled to the pooled cohort following NCHS guidance.

The LMF provides mortality follow-up time in months (\texttt{PERMTH\_EXM}) and the NCHS 113-cause underlying cause-of-death recode (\texttt{UCOD\_LEADING}), plus contributing-cause flags for diabetes and hypertension.

\subsection*{Participants and inclusion criteria}

Eligible participants were adults aged 40 to 79 years at NHANES examination with non-missing values on the metabolic syndrome components (waist circumference, fasting glucose, HDL cholesterol, triglycerides, and blood pressure). The age window was chosen to align with the PCE applicability range. Because triglycerides and HDL are measured only in the morning fasting subsample, the cohort is restricted to fasting-subsample participants, which removes approximately half of adult NHANES respondents. The HDL column name varies across cycles (LBDHDL in cycles A and B, 1999 to 2002; LBXHDD in cycle C, 2003 to 2004; LBDHDD in cycles D onward, 2005 to 2018) and was coalesced. Race and ethnicity were coalesced from RIDRETH3 (available from cycle G, 2011 onward) and RIDRETH1 (earlier cycles), with the limitation that the non-Hispanic Asian category is not separately identifiable before cycle G.

\subsection*{Risk scores}

Five risk scores were computed at NHANES baseline.

\emph{RMRS} \citep{shin2024rmrs} maps the five metabolic syndrome components to a continuous index via triangular areal similarity to a reference healthy point. Implementation follows the published formula without modification.

\emph{Metabolic syndrome decision tree} \citep{shin2023b9} classifies subjects into terminal nodes via a published CART. The published Figure 8(A) was transcribed; the left subtree (BPWC\_add $\leq 0.66$) is implemented as a single safety-zone leaf with the published probability of $p = 0.137$. A sensitivity analysis refitting the left subtree on a held-out NHANES split was planned (see Sensitivity analyses).

\emph{ACC/AHA Pooled Cohort Equations (PCE)} \citep{goff2014pce} estimate 10-year atherosclerotic cardiovascular disease risk. The four sex- and race-specific equations were implemented per the published coefficients. PCE applies to adults aged 40 to 79 without prior atherosclerotic cardiovascular disease; subjects outside this range or with missing race were excluded from PCE-specific analyses (PCE-eligible N reported in Results).

\emph{Framingham 2008 general cardiovascular risk profile} \citep{dagostino2008framingham} estimates 10-year general cardiovascular risk. Coefficients from the original publication were applied.

\emph{FINDRISC} \citep{lindstrom2003findrisc} is an integer screening score (range 0 to 26) for incident type 2 diabetes risk. NHANES does not collect FINDRISC items directly; the implementation maps NHANES variables onto FINDRISC items as follows. Age, BMI, and waist circumference map directly. Physical activity is approximated from the PAQ (physical activity questionnaire) items. Daily vegetable and fruit intake is approximated from DR1TOT (24-hour dietary recall). Antihypertensive medication use maps from BPQ050A. Personal history of high blood glucose maps from DIQ160. Family history of diabetes maps from MCQ300C. The proxy fidelity of these mappings is acknowledged as a study limitation; a sensitivity analysis with refined FINDRISC proxies was planned.

\subsection*{Outcomes}

Three mortality outcomes were defined.

\emph{Cardiovascular mortality} was defined as UCOD\_LEADING values for heart diseases (codes 1 to 3) or stroke (code 5). The follow-up cap was 10 years; AUC was reported at $t = 5$ and at the off-cap horizon $t = 9.5$ years (AUC at exactly the cap is degenerate under IPCW).

\emph{All-cause mortality} was defined as any death recorded in the LMF. The follow-up cap and horizons matched the cardiovascular outcome.

\emph{Diabetes-related mortality} used a broadened definition: UCOD\_LEADING for diabetes mellitus (code 7), OR UCOD\_LEADING for nephritis (code 9), OR the LMF DIABETES contributing-cause flag equal to 1. The follow-up cap was extended to 15 years to accumulate adequate events under the broadened definition; AUC horizons were $t = 5$, $t = 10$, and $t = 14.5$. The pre-registered narrow definition (UCOD\_LEADING $= 7$ only) yielded 6 events at 10 years and was underpowered for Fine-Gray modeling; the broadening to include nephritis and the contributing-cause flag, plus the follow-up extension, was registered as an amendment to the OSF protocol prior to running survival models on the diabetes outcome.

\subsection*{Statistical analysis}

For cardiovascular and diabetes-related mortality, we modeled cumulative incidence under competing risks using survey-weighted Fine-Gray subdistribution hazards \citep{finegray1999}. For all-cause mortality, we used survey-weighted Cox proportional hazards. Time-dependent discrimination was quantified with IPCW AUC \citep{blanche2013ipcwauc} via the \texttt{timeROC} R package with \texttt{iid = FALSE} (the \texttt{iid = TRUE} setting allocates an $N \times N$ matrix that exceeds available memory at the realized cohort size; see Software note in the Supplement).

Calibration was assessed via observed-vs-expected plots and the Greenwood-Nam-D'Agostino test. Net benefit was reported via decision curve analysis \citep{vickers2006dca} at clinically motivated threshold ranges (5\% to 20\% for all-cause, 5\% to 20\% for cardiovascular with the ACC/AHA 7.5\% statin threshold flagged, and 1\% to 3\% for diabetes-related). Because FINDRISC is an integer 0 to 26 score whereas the other four scores are probability-style, each score was Cox-recalibrated to a horizon-specific predicted probability before entering decision curve analysis. This per-score recalibration is reported with the net-benefit results and is necessary for cross-score comparability.

For the registered pairwise comparisons (H1: RMRS vs B9 on each outcome; H2: PCE vs Framingham on cardiovascular; H3: RMRS vs FINDRISC on diabetes-related), delta-AUC point estimates were taken from the IPCW \texttt{timeROC} fit; the corresponding 95\% CIs and DeLong p-values were computed via \texttt{pROC::roc.test} on the horizon-cap binary outcome (subjects censored before the horizon were excluded from the DeLong test). Definitive confidence intervals for all primary discrimination and net-benefit quantities were computed via PSU-cluster bootstrap at 500 resamples, sampling whole \texttt{sdmvstra}~$\times$~\texttt{sdmvpsu} clusters with replacement (about 301 clusters across pooled cycles). The bootstrap distribution is the registered sensitivity for the H3 non-inferiority test against the $-0.05$ AUC margin.

Incremental value of adding RMRS on top of FINDRISC for diabetes-related mortality was reported as continuous NRI \citep{pencina2008nri} and IDI, computed via \texttt{survIDINRI::IDI.INF} with 500 perturbation iterations.

Multiple imputation by chained equations was used for the small number of subjects with missing values on smoking, total cholesterol, or systolic blood pressure (PCE inputs). Twenty imputations were combined using Rubin's rules.

\subsection*{Sensitivity analyses}

The following sensitivity analyses were pre-registered (OSF section 8) and either completed or in progress.

\emph{ML upper-bound benchmark.} An XGBoost classifier trained on the five metabolic syndrome components plus age, sex, and race/ethnicity, with PSU $\times$ stratum-grouped 5-fold cross-validation, was used as a non-parametric ceiling for the achievable AUC from the same inputs. Hyperparameters were not tuned beyond defaults.

\emph{B9 left-subtree refit.} A CART was refit on a held-out NHANES split to test whether the original B9 left-subtree approximation contributes to the observed B9 underperformance. (Planned; status in \texttt{prereg/osf-preregistration-draft.md}.)

\emph{FINDRISC proxy refinement.} A second analysis with refined FINDRISC item proxies (improved mapping of PAQ, DR1TOT, and MCQ300C) was planned to bound the impact of proxy imprecision on the H3 result. (Planned.)

\emph{Subgroup analyses.} AUC and Fine-Gray models were planned within strata defined by sex, race/ethnicity, and age band (40 to 49, 50 to 59, 60 to 69, 70 to 79 years). The non-Hispanic Asian category is reported only for cycles G onward. (Planned.)

\emph{Bootstrap definitive CIs.} PSU-cluster bootstrap at 500 resamples for the focal AUCs and delta-AUCs is the registered definitive sensitivity for H3 and for all primary discrimination contrasts. (In progress at manuscript draft time.)

\subsection*{Reporting standards and pre-registration}

The study followed TRIPOD+AI \citep{tripodai2024} and STROBE \citep{strobe2007} reporting checklists; completed checklists are provided in the Supplement. The analysis plan was pre-registered on the Open Science Framework before any survival model was fit on the data; the registered protocol is available at the OSF identifier given in the Data and Code Availability section.

\subsection*{Software}

Analyses were conducted in R 4.3.3 on Ubuntu 24.04. Key packages included \texttt{survey} (NHANES design), \texttt{riskRegression} (Fine-Gray), \texttt{survival} and \texttt{rms} 6.7-1 (Cox), \texttt{timeROC} (IPCW AUC), \texttt{pROC} (DeLong delta-AUC), \texttt{dcurves} (DCA), \texttt{survIDINRI} (NRI and IDI), and \texttt{mice} (imputation). The XGBoost sensitivity used Python 3.11 with \texttt{xgboost} and \texttt{scikit-learn}. All code is available at the public repository listed in Data and Code Availability.

\section*{Results}

\subsection*{Cohort}

The pooled NHANES 1999 to 2018 fasting subsample after applying the inclusion criteria comprised 17{,}031 adults aged 40 to 79 years. Of these, 9{,}815 were within the PCE age and race applicability window and contributed to PCE-specific analyses. Within the 10-year follow-up cap, 806 all-cause deaths and 181 cardiovascular deaths accrued. Within the extended 15-year cap for diabetes-related mortality, 80 deaths met the broadened definition (UCOD\_LEADING for diabetes mellitus or nephritis, or the LMF DIABETES contributing-cause flag). A complete Table 1 of survey-weighted baseline characteristics by outcome status is provided in the Supplement.

\subsection*{Discrimination}

Time-dependent AUCs at the primary horizons are summarized in Table~\ref{tab:auc}. For all-cause mortality at the 9.5-year horizon, discrimination ranged from 0.525 (B9 decision tree) to 0.810 (Framingham 2008). For cardiovascular mortality at the same horizon, the range was 0.565 (B9) to 0.859 (Framingham). For diabetes-related mortality at the 14.5-year horizon, the range was 0.585 (B9) to 0.766 (FINDRISC), with RMRS at 0.747 sitting close to FINDRISC and well above B9.

\begin{table}[h]
\centering
\caption{Time-dependent AUC at primary horizons with 95\% PSU-cluster bootstrap CIs (500 resamples). PCE AUCs are computed on the PCE-eligible subset (N = 9{,}815). All other AUCs are computed on the full cohort (N = 17{,}031 for all-cause and cardiovascular; N = 16{,}996 for FINDRISC after listwise deletion on FINDRISC inputs).}
\label{tab:auc}
\begin{tabular}{llll}
\toprule
Score & Outcome & Horizon & AUC (95\% CI) \\
\midrule
Framingham 2008 & All-cause & 9.5y & 0.810 (0.795, 0.825) \\
PCE & All-cause & 9.5y & 0.758 (0.736, 0.778) \\
FINDRISC & All-cause & 9.5y & 0.682 (0.659, 0.703) \\
RMRS & All-cause & 9.5y & 0.595 (0.574, 0.618) \\
B9 tree & All-cause & 9.5y & 0.525 (0.505, 0.542) \\
\midrule
Framingham 2008 & Cardiovascular & 9.5y & 0.859 (0.833, 0.883) \\
PCE & Cardiovascular & 9.5y & 0.811 (0.776, 0.846) \\
RMRS & Cardiovascular & 9.5y & 0.662 (0.625, 0.692) \\
B9 tree & Cardiovascular & 9.5y & 0.565 (0.528, 0.597) \\
\midrule
FINDRISC & Diabetes-related & 14.5y & 0.766 (0.715, 0.817) \\
RMRS & Diabetes-related & 14.5y & 0.747 (0.692, 0.800) \\
B9 tree & Diabetes-related & 14.5y & 0.585 (0.533, 0.633) \\
\bottomrule
\end{tabular}
\end{table}

\subsection*{Calibration}

Observed-vs-expected calibration plots at each primary horizon are provided in the Supplement. The Greenwood-Nam-D'Agostino statistics indicated reasonable calibration for the clinical risk equations on their target outcomes and miscalibration for the metabolic-syndrome-derived scores when used as raw probabilities, which motivated the per-score Cox recalibration used for the decision curve analysis below.

\subsection*{Net benefit}

Decision curve analysis figures for the three outcomes are provided as \texttt{results/dca\_allcause.png}, \texttt{results/dca\_cv.png}, and \texttt{results/dca\_dm.png}. For all-cause mortality at the 9.5-year horizon and a 5\% decision threshold, PCE produced the largest net benefit (0.0331 per person, 95\% CI 0.0269 to 0.0390), followed by Framingham (0.0204, 0.0168 to 0.0241) and FINDRISC (0.0063, 0.0028 to 0.0100). RMRS and the B9 tree at this threshold sat at or below treat-none. For cardiovascular mortality at the ACC/AHA 7.5\% statin threshold, the net-benefit ordering followed the AUC ordering, with PCE and Framingham above the metabolic-syndrome-derived scores at all clinically relevant thresholds in the 5\% to 20\% range. For diabetes-related mortality in the 1\% to 3\% threshold range, FINDRISC and RMRS exceeded treat-none across the range while B9 did not.

\subsection*{Pairwise comparisons and incremental value}

The five registered pairwise delta-AUC comparisons are summarized in Table~\ref{tab:pairwise}.

\begin{table}[h]
\centering
\caption{Registered pairwise delta-AUC comparisons. CIs are 95\% PSU-cluster bootstrap percentile intervals (500 resamples); DeLong-based CIs agreed to two decimals across all five contrasts.}
\label{tab:pairwise}
\begin{tabular}{lll}
\toprule
Contrast & Outcome (horizon) & Delta-AUC (95\% CI) \\
\midrule
RMRS vs FINDRISC (H3, non-inferiority) & Diabetes-related (14.5y) & $+0.011$ ($-0.053$, $+0.075$) \\
RMRS vs B9 & All-cause (9.5y) & $+0.066$ ($+0.044$, $+0.089$) \\
RMRS vs B9 & Cardiovascular (9.5y) & $+0.086$ ($+0.044$, $+0.129$) \\
RMRS vs B9 & Diabetes-related (14.5y) & $+0.149$ ($+0.079$, $+0.215$) \\
PCE vs Framingham (H2) & Cardiovascular (9.5y) & $+0.009$ ($-0.008$, $+0.026$) \\
\bottomrule
\end{tabular}
\end{table}

For the H3 non-inferiority test of RMRS versus FINDRISC on diabetes-related mortality at the 14.5-year horizon, the bootstrap 95\% percentile lower bound was $-0.053$, which exceeded the registered margin of $-0.05$ by 0.003 AUC units. The point estimate of $+0.011$ favored RMRS. The DeLong-based 95\% CI on the horizon-cap binary outcome was $(-0.055, +0.077)$, agreeing with the bootstrap to a thousandth. By the registered decision rule, non-inferiority was not formally established.

Adding RMRS on top of FINDRISC for diabetes-related mortality produced an IDI of 0.0031 (95\% CI 0.0002 to 0.0099) and a continuous NRI of 0.173 (95\% CI 0.033 to 0.287). Both intervals excluded zero, indicating that RMRS captured diabetes-mortality information beyond what FINDRISC already provided, with modest magnitude.

For the H2 comparison of PCE versus Framingham on cardiovascular mortality, the delta-AUC was $+0.009$ (95\% CI $-0.008$ to $+0.026$) favoring PCE. The interval contained zero, indicating no detectable difference between the two clinical equations on this outcome in this cohort.

\subsection*{Sensitivity analyses}

The XGBoost ceiling on the five metabolic syndrome components plus age, sex, and race/ethnicity, with PSU $\times$ stratum-grouped 5-fold cross-validation, yielded AUCs of 0.813 for all-cause mortality at 9.5 years, 0.813 for cardiovascular mortality at 9.5 years, and 0.799 for diabetes-related mortality at 14.5 years. For all-cause mortality, the ML ceiling coincided with the Framingham AUC, indicating no headroom for a tree-based metabolic syndrome variant on this outcome. For cardiovascular mortality, the clinical Framingham equation at 0.859 exceeded the ML ceiling at 0.813, indicating that the metabolic syndrome inputs alone cannot reach the Framingham performance level even under flexible non-parametric modeling. For diabetes-related mortality, the ML ceiling at 0.799 exceeded RMRS at 0.747 and FINDRISC at 0.766, leaving roughly 5 AUC points of headroom for a tree-based variant of the metabolic syndrome representation specifically for diabetes mortality.

The B9 left-subtree US-refit sensitivity, the FINDRISC proxy refinement, and the pre-registered subgroup analyses by sex, race/ethnicity, and age band are provided in the Supplement (status noted at submission time).

\section*{Discussion}

In this external longitudinal validation of two metabolic syndrome risk scoring methods (RMRS and the B9 decision tree) against three established clinical risk equations (PCE, Framingham 2008, FINDRISC) on the US NHANES Linked Mortality File, the headline finding is outcome-specific. The clinical risk equations dominated the metabolic-syndrome-derived scores when the outcome was cardiovascular or all-cause mortality. For diabetes-related mortality, the metabolic-syndrome-derived RMRS competed with FINDRISC on point discrimination and substantially outperformed the B9 decision tree.

For diabetes-related mortality at the 14.5-year horizon, RMRS achieved an AUC of 0.747 against FINDRISC at 0.766, with a delta-AUC of $+0.011$ favoring RMRS and a bootstrap 95\% percentile CI of $(-0.053, +0.075)$. The registered non-inferiority test at the $-0.05$ margin was not formally established by either the bootstrap or the DeLong interval, although both intervals locate the two scores within a 5 AUC point band on this outcome. The IDI and continuous NRI for adding RMRS on top of FINDRISC excluded zero, indicating that RMRS captures diabetes-mortality information beyond the FINDRISC items, with magnitude on the order of one to three percentage points of reclassification. RMRS occupies a complementary niche in diabetes-mortality prediction alongside FINDRISC, which remains the primary screening tool of choice.

For cardiovascular mortality at the 9.5-year horizon, Framingham 2008 achieved an AUC of 0.859 and PCE achieved 0.811, with a delta-AUC of $+0.009$ between them (95\% CI $-0.008$ to $+0.026$). The metabolic-syndrome-derived RMRS at 0.662 sat about 20 AUC points below these clinical equations. The decision curve analysis at the ACC/AHA 7.5\% statin threshold showed the same ordering. For all-cause mortality, Framingham again led at 0.810, with PCE at 0.758 and the metabolic-syndrome-derived scores in the 0.525 to 0.595 range. The XGBoost ceiling on the metabolic syndrome inputs alone, at 0.813 for both cardiovascular and all-cause, did not close the gap. The metabolic syndrome features do not contain the cardiovascular-mortality and all-cause-mortality information that the clinical equations integrate, regardless of the modeling approach used on them.

The B9 decision tree formulation consistently underperformed the continuous RMRS on every outcome, with delta-AUCs against RMRS of $+0.066$ for all-cause, $+0.086$ for cardiovascular, and $+0.149$ for diabetes-related mortality (all CIs excluded zero). The XGBoost ceiling, fit on the same five metabolic syndrome inputs, exceeded both the continuous RMRS and the B9 tree on diabetes mortality (0.799 vs 0.747 vs 0.585). Taken together, these patterns suggest that the B9 underperformance is attributable to the tree-structure decision and not to the metabolic syndrome signal itself. A US-population CART refit on a held-out NHANES split, listed as a pre-registered sensitivity in the Supplement, will further disambiguate the contribution of transportability and methodology to the gap.

Several limitations apply. First, the FINDRISC implementation in NHANES relies on proxy mappings of physical activity, daily vegetable and fruit intake, and family history of diabetes onto NHANES PAQ, DR1TOT, and MCQ300C items respectively. The proxy fidelity is imperfect; a sensitivity analysis with refined mappings is reported in the Supplement. Second, the cohort is restricted to the morning fasting subsample of NHANES, which removes approximately half of adult NHANES respondents and may introduce selection on factors associated with willingness or ability to fast. Third, the non-Hispanic Asian race category is only separately identifiable in NHANES cycles G onward (2011 to 2018), which limits subgroup precision for that group. Fourth, the diabetes-related mortality outcome uses a broadened definition (UCOD for diabetes mellitus or nephritis, or the LMF diabetes contributing-cause flag); the pre-registered narrow definition was underpowered for Fine-Gray modeling with only 6 events at the 10-year horizon, and the broadening was registered as an amendment prior to running survival models on this outcome. The clinical implication is that the metabolic syndrome continuous risk score occupies a complementary niche in diabetes-mortality prediction (alongside FINDRISC), and does not substitute for established clinical risk equations for cardiovascular or all-cause mortality prediction. Future work should evaluate the metabolic syndrome inputs combined with the clinical risk equation inputs in joint models, and validate the diabetes-mortality finding in non-US cohorts.

\section*{Data and code availability}

All analysis code is available at \url{https://github.com/jayhemnani9910/longitudinal-mets-validation} (MIT license). The OSF pre-registration is at (TBD: OSF identifier to be inserted on submission).

\section*{Conflicts of interest}

The author reports no conflicts of interest.

\section*{Acknowledgments}

The author thanks Shin, Oh, and Shim for the underlying RMRS and B9 publications that motivated this external validation.

\bibliographystyle{plainnat}
\bibliography{main}

\end{document}
